\section{Full Procurement Runtime Specification}
\label{app:runtime}

This appendix makes the main, two-stage procurement system explicit before describing the
single-call algebra used by PACS and WTQ.  The two representations share a deterministic backend,
but they are not interchangeable and are never combined within one reported prediction.

\subsection{Intermediate representations}

Step 1 returns an \emph{intent program}.  Its answer signature declares the operation, value type,
and requested answer field; its ordered steps declare literal filters, dependencies, and the final
answer step.  An unsupported or ambiguous item instead carries a reason and an empty official
program.  A valid intent is either compiled directly or supplied as a briefing to Step 2.

Step 2 returns the graph-plan interface in Table \ref{tab:graph-ir}.  This representation is used
for a multi-hop dependency such as records $\rightarrow$ emitted suppliers $\rightarrow$ records
bound by those suppliers $\rightarrow$ a scalar or set answer.  The compiler normalises supported
planner idioms, constructs a directed acyclic graph, and lowers every executable variable to a
backend query specification.

\begin{table}[h]
  \centering
  \caption{Graph-plan intermediate representation used by the full procurement runtime.}
  \label{tab:graph-ir}
  \small
  \begin{tabularx}{\linewidth}{p{0.19\linewidth}p{0.25\linewidth}X}
    \toprule
    Object & Required content & Deterministic checks \\
    \midrule
    Variable & identifier, kind, semantic role, literal constraints, optional emitted field & unique identifier; legal field and value type; no invented root input \\
    Dependency & upstream variable and downstream consumer & reference exists; topological order; no cycle \\
    Relation & source variable/field and target variable/bind field & emitted and bound types agree; no unconstrained bind \\
    Return & input variable, operation, answer field/type, optional grouping or comparator & operation accepts the variable kind; requested count target and direction agree with the question \\
    Guards & exhaustive reduction, deduplication key, additivity and optional population conditions & safe aggregate and declared backend semantics \\
    Abstention & unsupported or ambiguous reason & no answering graph may coexist with a terminal non-answer decision \\
    \bottomrule
  \end{tabularx}
\end{table}

\subsection{Online algorithm}

For one question, the implemented runtime performs the following steps.

\begin{enumerate}
  \item Select a schema slice from question keywords.  This is rule-based schema selection, not
  dense retrieval over records.
  \item Ask the Step-1 model for an intent program at temperature zero.  Validate references,
  dependency order, return type, answer signature, and any terminal non-answer reason.
  \item If the intent is valid, compilable, and not vetoed, compile it deterministically.  Otherwise
  call the Step-2 typed graph planner.  In the evaluation protocol, at most two structural samples
  are considered, and a sample is retried only after a detectable compile or consistency failure.
  \item Ground aliases, fields, entity literals, emitted values, downstream bindings, aggregation
  guards, and the deduplication key.  Run schema, conflict, and exhaustive-reduction preflight.
  \item Execute graph variables in dependency order.  Derived value sets become hidden membership
  constraints on downstream record variables.  Compute the declared return operation from the
  resulting record population.
  \item Run postflight, answer-sanity, and evidence-release checks.  Return an answer card only when
  the resulting state is releasable.  A clean no-result, unsupported, or ambiguous state is final.
  \item If a compile, grounding, or execution defect is explicitly repairable, supply the structured
  failure to the planner and repeat Steps 3--6.  Main evaluation permits one feedback replan;
  teacher collection permits two.
\end{enumerate}

The online path never receives a gold answer.  Structural sampling therefore cannot repair an
executable-but-semantically-wrong plan unless a runtime check detects its defect.

The current compiler also has a known fidelity gap.  Its post-compile check can add a missing year,
CPV, or category atom to the flat query specification, but the pipeline executes an existing graph
plan in preference to that flat specification.  Until the graph is recompiled or its terminal
variable is updated, this completion is not a guarantee that the executed graph contains the atom.

\subsection{Evidence and release contract}

Table \ref{tab:evidence-contract} distinguishes exact execution state from bounded audit material.
This distinction prevents a sample of identifiers from being mistaken for the complete population
used by a count or sum.

\begin{table}[h]
  \centering
  \caption{Information produced and retained by the procurement runtime.}
  \label{tab:evidence-contract}
  \small
  \begin{tabularx}{\linewidth}{p{0.24\linewidth}p{0.28\linewidth}X}
    \toprule
    Quantity & Runtime semantics & Retention boundary \\
    \midrule
    Answer and population count & exact deterministic result under the frozen record view & stored in a rich trace; prediction and correctness are retained by the slim main evaluator \\
    Variable trace & grounded constraints, dependencies, emitted field, status, and output size & available when rich tracing is enabled; omitted from the main 2,285-row result files \\
    Contract evidence & matching contract identifiers and rows & exact total count, but identifiers/rows are capped at 200 per execution result \\
    Public evidence view & operation, answer, count, identifiers, and OCIDs & serialises at most 50 identifiers and 50 OCIDs from the runtime bundle \\
    Source documents & URL and document metadata extracted from OCDS & no production document inspector or passage-level citation scorer \\
    Oracle & typed answer used for offline curation and scoring & never included in online prompts or student targets \\
    \bottomrule
  \end{tabularx}
\end{table}

\subsection{Illustrative trace provenance}

The candidate running example supplied for Figure \ref{fig:overview} is
\texttt{L2::bridge\_join\_0502\#L2a}: ``Record check: how many contract notices were published by
buyers that have awarded a contract to Dodd Group (Midlands) Limited?''  The saved Step-1 intent takes the
deterministic compiler branch; Step 2 is bypassed.  A no-model replay on the current frozen backend
finds 21 source records, emits 16 buyer values, binds those values into 1,940 target records, and
returns 1,940, matching the stored oracle.  The release includes a limitation: eight of the first
200 matched rows have no recorded supplier field.  This is a reconstruction from a saved Step-1
artifact plus current deterministic replay, not a fresh final-test model call.  The exact source
locations and counts are recorded in
\path{paper/tmlr_v3/figs/src/fig01_runtime_bridge.manifest.json}; a replacement hand-drawn figure
should preserve this qualification or use a new rich trace with its own manifest.

\section{Evaluation Protocol Matrix}
\label{app:protocol-matrix}

Table \ref{tab:protocol-matrix} records the decoder and authority boundaries needed to interpret
the four result sets.  A dash means that the component does not exist in that protocol, rather than
that it failed to run.

\begin{table}[h]
  \centering
  \caption{Inference and curation protocols.  ``Oracle'' refers to access during the named
  procedure, not to whether an oracle later exists for scoring.}
  \label{tab:protocol-matrix}
  \small
  \begin{tabularx}{\linewidth}{p{0.17\linewidth}p{0.19\linewidth}p{0.17\linewidth}p{0.13\linewidth}X}
    \toprule
    Protocol & Plan form & Decoding & Repair & Oracle and saved output \\
    \midrule
    Teacher harvest (9,267) & Step-1 intent; conditional Step-2 graph & temperature 0; two structural samples & at most 2 & oracle-hidden online, then oracle/shape gate offline; rich trace and separate training sinks \\
    Procurement main (2,285; primary 2,025) & same full runtime & temperature 0; two structural samples & at most 1 & no online oracle; saved evaluator files retain slim per-item predictions \\
    Development (260) & same full runtime; checkpoint varies & temperature 0; two structural samples & at most 1 & no online oracle; paired predictions by question \\
    PACS (922) & one recursive algebra tree & temperature 0; primary protocol is prompt-only JSON & none & no online oracle; tree, execution status, answer, and offline score saved \\
    WTQ (4,344) & one recursive algebra tree with per-table schema & temperature 0; JSON-schema guided & none in the reported pristine run & no online oracle; answer items rescored by the official evaluator \\
    \bottomrule
  \end{tabularx}
\end{table}

The main procurement run directories do not preserve one complete command and checkpoint manifest.
The PACS corrected per-item file and the exact WTQ C-final training snapshot are also absent.  These
are historical artifact gaps, not omitted algorithmic choices; they are stated here to delimit the
level of reproduction currently possible.

\section{Recursive Algebra}
\label{app:algebra}

The direct-composition language is a closed tree algebra.  Programs are limited to depth 16 and
64 nodes.  Table \ref{tab:algebra} gives the public node signatures; predicates include equality,
membership, substring containment, existence, numeric bounds, negation, disjunction, and
membership in a computed \values{} set.

\begin{table}[h]
  \centering
  \caption{Recursive-algebra operators.  The validator checks compositional types and required
  arguments before execution.}
  \label{tab:algebra}
  \small
  \begin{tabular}{lll}
    \toprule
    Node & Signature & Function \\
    \midrule
    \texttt{filter} & predicates $\rightarrow$ \records & filter records \\
    \texttt{values} & \records $\rightarrow$ \values & project distinct values \\
    \texttt{count}, \texttt{sum} & \records $\rightarrow$ \numbertype & aggregate records \\
    \texttt{size} & \values $\rightarrow$ \numbertype & set cardinality \\
    \texttt{exists} & \records $\rightarrow$ \booltype & test non-emptiness \\
    \texttt{select}, \texttt{extreme} & \records $\rightarrow$ \valuetype & scalar or extremal lookup \\
    \texttt{groupby} & \records $\rightarrow$ \groups & grouped count or sum \\
    \texttt{argext}, \texttt{top} & \groups $\rightarrow$ scalar/ranking & rank groups \\
    \texttt{combine} & \numbertype$^2\rightarrow$ scalar & arithmetic or comparison \\
    \texttt{vcompare} & \valuetype$^2\rightarrow$ \booltype & compare values \\
    \texttt{setop} & \values$^2\rightarrow$ \values & union, intersection, difference \\
    \texttt{gcombine} & \groups$^2\rightarrow$ \groups & combine aligned groups \\
    \texttt{keys\_where} & \groups$\rightarrow$ \values & threshold group keys \\
    \bottomrule
  \end{tabular}
\end{table}

Counts and grouped counts deduplicate contract identifiers.  Projections remove null and empty
values and return sorted distinct strings.  Sums use decimal accumulation.  Ranking ties are
ordered by their stored string keys.  A Boolean matches only a Boolean, so \texttt{True} is not
equal to numeric one.  The type checker does not by itself enforce every domain convention:
additivity, schema membership, and currency remain separate contracts.

\section{Data Construction and Schema}
\label{app:data}

\paragraph{Release selection.}
The loader streams yearly compressed JSON Lines files, skips malformed rows, flattens the required
OCDS fields, and retains the latest parseable release per OCID.  It also stores the years in which
the process was observed.  This creates a current snapshot rather than a merged event history.

\paragraph{Entity resolution.}
Registry identifiers have priority.  Platform identifiers are linked only under the same OCID,
normalised name, and role, and only if all observations identify one official entity.  The second
phase uses government lookup, exact name and region, then exact name.  Empty normalised names remain
singletons.  Fuzzy similarity creates review candidates only.  Reviewed safe edges are applied by
a cluster consistency check.  The current artifacts contain 204,711 raw aliases and 131,502
canonical organisations; 82,013 remain singletons.

\paragraph{Values and provenance.}
The amount source precedence is award, linked contract, then tender.  Award- and contract-sourced
values may be labelled additive; tender ceilings are not.  Numeric zero is retained.  Currency is
stored but not converted.  Evidence nodes retain the source field and source document linkage.

\begin{table}[h]
  \centering
  \caption{Retained graph snapshot.  One contract may have several supplier edges.}
  \label{tab:kg-schema}
  \small
  \begin{tabular}{lr@{\hskip 3em}lr}
    \toprule
    Node type & Count & Edge type & Count \\
    \midrule
    Contract & 215,221 & \texttt{buyer\_of} & 215,218 \\
    Organisation & 131,502 & \texttt{supplier\_of} & 334,063 \\
    CPV & 3,870 & \texttt{categorized\_by} & 164,691 \\
    Evidence & 535,731 & \texttt{evidence\_for} & 1,326,240 \\
    \bottomrule
  \end{tabular}
\end{table}

\section{Benchmark Construction and Integrity}
\label{app:benchmark}

The initial procurement generator creates executable answer specifications from the graph.  Gate A
checks conflicts, identifier integrity, deterministic execution, and aggregate sanity before
surface generation.  A second language-model stage writes the question; Gate B checks that visible
filters and operations agree with the specification.  The initial artifacts contain 10,000 answer
specifications, 9,053 generated questions, and 9 Gate-B failures, yielding 9,044 rows.  Later
curation adds answerability twins, whole-plan split assignment, class balancing, and surface
variation to obtain 12,828 graph-benchmark rows: 9,267 train, 671 development-select, 556
development-tune, 49 smoke, and 2,285 final-test rows.

Plan identifiers do not cross from train to final test in the stored split.  A separate string
audit nevertheless finds two identical question--answer pairs and five deduplication groups across
that boundary.  The split should therefore be described as plan-ID separated, not fully
duplicate-free.

\begin{table}[h]
  \centering
  \caption{Composition of the 2,285-row frozen procurement partition.  The 260-row development
  subset contains 20 rows from each bucket and is removed by identifier to form the 2,025-row
  primary column.}
  \label{tab:final-buckets}
  \small
  \begin{tabular}{lr@{\hskip 2.2em}lr@{\hskip 2.2em}lr}
    \toprule
    Bucket & Rows & Bucket & Rows & Bucket & Rows \\
    \midrule
    Ambiguous & 120 & Boolean & 121 & Bridge join & 350 \\
    No result & 120 & Categorical & 120 & Comparison & 306 \\
    Unsupported & 120 & Count & 250 & Factoid & 250 \\
    Min/max & 107 & Set & 93 & Sum & 200 \\
    Top-$k$ & 128 & Answerable total & 1,925 & Grand total & 2,285 \\
    \bottomrule
  \end{tabular}
\end{table}

\begin{table}[h]
  \centering
  \caption{Capability breakdown on the 260-question development diagnostic (20 items per row).
  Exact counts are reported because the balanced cells are small.}
  \label{tab:bucket-analysis}
  \small
  \begin{tabular}{lrrr@{\qquad}lrrr}
    \toprule
    Capability & KG-only & RAG-40 & Local Q & Capability & KG-only & RAG-40 & Local Q \\
    \midrule
    Ambiguous & 15 & 5 & 19 & Comparison & 8 & 3 & 17 \\
    No result & 20 & 13 & 20 & Count & 17 & 3 & 20 \\
    Unsupported & 20 & 20 & 19 & Factoid & 5 & 17 & 15 \\
    Boolean & 11 & 0 & 14 & Min/max & 18 & 0 & 20 \\
    Bridge join & 7 & 0 & 14 & Set & 9 & 1 & 15 \\
    Categorical & 5 & 13 & 13 & Sum & 19 & 0 & 20 \\
    & & & & Top-$k$ & 1 & 0 & 18 \\
    \bottomrule
  \end{tabular}
\end{table}

\section{PACS Details}
\label{app:pacs}

PACS instantiates seven families at levels L1--L3.  Construction samples executable typed trees,
runs the production and separate evaluators, requires answer agreement and non-degeneracy, and
excludes exact tree hashes found in Compose-v3 training or validation.  Missing-operator and
no-result controls are added separately.  The evaluated surface files contain 922 rows.  Table
\ref{tab:pacs-composition} reports the actual status and exposure labels used by evaluation.

\begin{table}[h]
  \centering
  \caption{Corrected PACS v1.1 composition.  ``Seen'' refers to the declared program-shape label,
  not exact tree overlap.}
  \label{tab:pacs-composition}
  \small
  \begin{tabular}{lr@{\hskip 3em}lr}
    \toprule
    Status & Rows & Shape label & Rows \\
    \midrule
    Answerable & 730 & Seen & 431 \\
    Unsupported & 138 & Unseen & 491 \\
    No result & 54 & Total & 922 \\
    \bottomrule
  \end{tabular}
\end{table}

The original surface channels treated 36 answerable false or zero-valued denotations as no-result.
The corrected v1.1 manifest restores them to the answerable class.  Saved predictions were joined
by item identifier and their trees were re-executed; no model was called.  This changes Compose-v3
channel A from 687 to 722 correct and channel B from 675 to 711 correct.  The repository does not
retain the original status-transformation script, so the paper artifact records this deterministic
correction explicitly.  The seal implementation also differs from its comments.  It checks exact
question and tree hashes, entity anchors, and a sampled trigram screen, but it does not implement a
separate template-identity gate and compares only every seventh training trigram set.  These facts
motivate the qualified description in Section \ref{sec:experiments}.

The declared seen/unseen labels mix answerable and status rows, so they are a difficulty diagnostic
rather than a pure measure of logical-form novelty.  Exact question and tree-hash overlap with the
PACS test is absent; this is a narrower claim than saying that every program shape is unseen.

\section{WikiTableQuestions Evaluation}
\label{app:wtq}

The WTQ evaluation uses 4,344 questions over 421 test tables absent from the training-table
partition.  Each table is exposed through two aligned views.  The typed view adds deterministic
numeric, date, year, month, and normalised-string columns for computation; the raw view preserves
the strings returned to the official scorer.  Row identifiers and derived columns are never shown
as answer values.

For official scoring, each predicted answer item first replaces tab and newline characters with a
space.  The question identifier and cleaned items are then joined with tab delimiters and passed to
\path{scripts/wtq/official_evaluator.py} version 1.0.2 with the original CoreNLP-tagged targets.
This detail changes three borderline items relative to calling the normaliser directly.  Table
\ref{tab:wtq-official} records the official counts and paired transitions reconstructed from saved
answer items.

\begin{table}[h]
  \centering
  \caption{WTQ pristine-unseen official evaluation.  Gain/loss compares each row with the row
  immediately above using exact McNemar tests.}
  \label{tab:wtq-official}
  \small
  \begin{tabular}{lrrl}
    \toprule
    Planner or supervision & Correct / 4,344 & Accuracy (\%) & Gain/loss; $p$ \\
    \midrule
    Untuned Qwen3-8B & 978 & 22.51 & -- \\
    Procurement Compose-v3 & 1,187 & 27.33 & 387/178; $8.50\times10^{-19}$ \\
    Answer-only self-harvest (A) & 1,930 & 44.43 & 833/90; $1.77\times10^{-151}$ \\
    Translated gold programs (C) & 2,250 & 51.80 & 628/308; $6.20\times10^{-26}$ \\
    \bottomrule
  \end{tabular}
\end{table}

The A run records 4,709 answer-matched examples.  The C run used approximately 5,423 translated
program examples selected by translation, type, execution, and answer agreement.  Its exact
training snapshot is not present in the current repository: the historical commit, current file,
and staged file contain different later reconstructions.  The official per-item predictions remain
available, but this missing snapshot prevents a turnkey retraining claim.

\section{Grounding and Target Audits}
\label{app:grounding-audit}

The fixed-plan replay executes the same saved proposal before and after deterministic grounding on
the same frozen backend.  Of 136 plans with typed oracles, 89 are correct before grounding and 103
after it.  Fourteen move from wrong to correct, none moves from correct to wrong, and one previously
non-executable plan becomes executable but remains incorrect.  All 14 rescues are sums for which
grounding inserts the additive-value guard.  This isolates one runtime transformation; it does not
estimate the effect of grounding-derived rows on a trained model.

The teacher artifacts contain 237 successful repairs attributed to schema, runtime-specification,
or entity-grounding failures.  Of 229 targets that can be checked mechanically against mapped
question constraints, 173 contain every mapped constraint and 56 are flagged.  Category is absent
from 43 flagged mappings, followed by buyer (6), supplier (3), signed-date guard (3), title (2), and
CPV (2); categories may overlap within one target.  A missing condition can be extensionally
redundant on one graph snapshot, allowing an incomplete plan to match the answer oracle.  This is
why the paper does not equate denotation agreement with full semantic equivalence.

\section{Training Data and Configurations}
\label{app:config}

Table \ref{tab:training-data} separates candidate harvest counts from final exported training
splits.  A candidate count is not an accuracy and overlapping objectives do not create new unique
questions.

\begin{table}[h]
  \centering
  \caption{Graph-plan teacher harvest and capped exports.}
  \label{tab:training-data}
  \small
  \begin{tabularx}{\linewidth}{XrX}
    \toprule
    Artifact & Harvest rows & Exported train/validation \\
    \midrule
    Training questions & 9,267 & -- \\
    Traces producing an answer after runtime checks & 6,860 & -- \\
    Oracle matches before shape filtering & 5,605 & -- \\
    Direct graph target pool passing both offline gates & 5,598 & included subject to caps \\
    Hard negative routes after an offline gate failure & 1,262 & -- \\
    Successful repair targets & 1,725 & 1,679 / 46 \\
    Preference pairs & 390 & 390 / -- \\
    Abstention target pool & 590 & included subject to a ten percent cap \\
    Final plan SFT export & -- & 2,787 / 57 \\
    Step-1 intent targets & -- & 5,091 / 96 \\
    \bottomrule
  \end{tabularx}
\end{table}

Table \ref{tab:training-config} records the configurations actually consumed by the reported
checkpoints.  Every row uses 4-bit QLoRA on all target modules, rank 64, $\alpha=128$, dropout
0.05, bfloat16 compute, effective batch 16, cosine decay, and one CUDA process.  Step-2 SFT trains
on 4,466 rows comprising 2,787 plan SFT rows and 1,679 repair targets.  Its configured validation set contains the
57 plan SFT rows; the separately exported 46 repair validation rows are not selected by the
training YAML.

\begin{table}[h]
  \centering
  \caption{Training and continuation configurations.  Q/L denotes matched Qwen3-8B and
  Llama-3.1-8B variants.}
  \label{tab:training-config}
  \small
  \begin{tabularx}{\linewidth}{p{0.15\linewidth}p{0.23\linewidth}p{0.17\linewidth}p{0.10\linewidth}p{0.10\linewidth}X}
    \toprule
    Stage & Initialisation and train rows & Objective & Epochs & LR & Cutoff / validation \\
    \midrule
    Step-2 SFT (Q/L) & base; 2,787 plan + 1,679 repair & token SFT & 3 & $10^{-4}$ & 6,144 / 57 plan \\
    Step-1 SFT (Q/L) & base; 5,091 intent programs & token SFT & 3 & $10^{-4}$ & 6,144 / 96 \\
    Qwen RSFT & Qwen SFT; 3,019 plan + 3,169 repair & token SFT & 2 & $5\times10^{-5}$ & 6,144 / 65 plan \\
    Llama RSFT & Llama SFT; 3,026 plan + 3,190 repair & token SFT & 2 & $5\times10^{-5}$ & 6,144 / 66 plan \\
    Qwen DPO & Qwen RSFT; 390 teacher + 689 on-policy pairs & sigmoid DPO, $\beta=.1$ & 1 & $5\times10^{-6}$ & 6,144 / none \\
    Llama DPO & Llama RSFT; 390 teacher + 693 on-policy pairs & sigmoid DPO, $\beta=.1$ & 1 & $5\times10^{-6}$ & 6,144 / none \\
    Compose-v3 & fresh Qwen; 12,414 trees & token SFT & 2 & $10^{-4}$ & 4,096 / 253 \\
    WTQ A & fresh Qwen; 4,615 answer-matched trees & token SFT & 2 & $10^{-4}$ & 4,096 / 94 \\
    WTQ C & fresh Qwen; about 5,423 translated programs & token SFT & 2 & $10^{-4}$ & 4,096 / snapshot missing \\
    \bottomrule
  \end{tabularx}
\end{table}

Training logs record Transformers 5.6.0 and the expected example and optimiser-step counts.  The
evaluation server logs record vLLM 0.24.0, seed 0, bfloat16, maximum sequence length 8,192, LoRA
rank 64, and disabled Qwen thinking.  They do not record a formal GPU-model manifest.  The paper
therefore reports one-device execution without claiming a hardware model that cannot be recovered
from the preserved logs.

\section{Artifact and Reproduction Notes}
\label{app:artifacts}

The implementation entry points are:
\begin{itemize}
  \item \path{pipelines/01_ingest.py} through \path{pipelines/04_extract.py}, followed by
  \path{pipelines/40_build_kg.py} and \path{pipelines/41_validate_kg.py};
  \item \path{scripts/apply_safe_er_merges.py} for the reviewed cluster-safe ER layer;
  \item \path{scripts/run_teacher.py} and \path{scripts/export_llamafactory.py} for graph-plan
  curation and export;
  \item \path{scripts/run_compare.py} for the full runtime, and
  \path{scripts/run_compose_probe_eval.py} for direct composition;
  \item \path{scripts/analyze_grounding_impact.py} for the fixed-plan replay.
\end{itemize}

Paper-side derived counts and provenance notes are stored in
\path{paper/tmlr_v3/artifacts/derived_metrics.json} and
\path{paper/tmlr_v3/artifacts/ARTIFACT_LEDGER.md}.  The compact KG-only development summary records
the SHA256 values of the 260-row input, runner, full per-item result, and source summary without
copying the 9 MB diagnostic result into the paper directory.

The prompt/call contract, source and configuration hashes, selection authority, serving settings,
and ablation classification are recorded in
\path{paper/tmlr_v3/artifacts/REPRODUCIBILITY_MANIFEST.md}.  The historical teacher summary stores
the Step-1 and Step-2 model names but omits the command, prompt hash, worker count, and most call
parameters.  Those fields are explicitly marked as reconstructed from the audited runner.  The
main procurement run directories likewise lack a unified arm-to-checkpoint manifest, and the
currently linked adapter directories are not portable repository artifacts.

Full production builds use canonical output paths.  Partial year, limit, or sample runs require an
explicit non-canonical output location, preventing diagnostic runs from replacing the canonical
tables.  The current working tree contains staged data artifacts and unstaged correctness fixes;
a release version must record a source commit, data hashes, checkpoint hashes, decoder settings,
and scorer version together.  We bind the present paper to its enumerated artifacts rather than
claiming that every historical summary can be regenerated from an unspecified working tree.

\section{Project-Brief Traceability}
\label{app:brief}

Table \ref{tab:brief-coverage} maps the original project requirements to implemented evidence and
remaining gaps.  It distinguishes a completed component from one that is designed or only partly
evaluated.

\begin{table}[h]
  \centering
  \caption{Coverage of the original LLM--KG procurement project brief.}
  \label{tab:brief-coverage}
  \small
  \begin{tabularx}{\linewidth}{p{0.25\linewidth}p{0.14\linewidth}X}
    \toprule
    Requirement & Status & Evidence and boundary \\
    \midrule
    Structured/semi-structured ETL and KG schema & Implemented & versioned OCDS ingestion, rich extraction, four node types, four edge types, validation before write \\
    Entity resolution and aliases & Implemented & registry-first and precision-first deterministic rules, reviewed safe merges, alias audit, unresolved singletons \\
    Natural language to executable KG computation & Implemented & typed Step-1 intent, deterministic fast path, conditional graph planner, grounding and executor \\
    Multi-hop, aggregation, temporal and set reasoning & Implemented & bridge dependencies, count/sum/extrema/top-$k$/set/Boolean operations, year and date constraints \\
    Evidence-linked reasoning trace & Partial & variable traces and contract identifiers exist; main evaluation stores slim outputs and no source-passage citation score \\
    Reliability controls & Implemented & typed/schema checks, exhaustive execution, pre/postflight, answer sanity, explicit abstention, bounded replan \\
    Factual, compositional and non-answer evaluation & Implemented & procurement final partition, PACS, and explicit ambiguous/unsupported/no-result classes \\
    Citation accuracy & Not evaluated & document metadata exist, but there is no production passage inspector or citation gold set \\
    Sensitivity to KG incompleteness & Not evaluated & no controlled node/edge deletion or missing-field intervention study \\
    Closed-book, RAG, LLM+KG, and deterministic controls & Implemented & all four families are reported; the deterministic control is rule/decomposition over the KG rather than SPARQL/Cypher \\
    \bottomrule
  \end{tabularx}
\end{table}

\section{Ablation Coverage and Causal Scope}
\label{app:ablations}

Table \ref{tab:ablation-scope} separates a controlled intervention from a checkpoint or robustness
comparison.  This distinction matters because a fixed test set alone does not make a training
ladder a component ablation.

\begin{table}[h]
  \centering
  \caption{What each existing comparison holds fixed and what it can establish.}
  \label{tab:ablation-scope}
  \small
  \begin{tabularx}{\linewidth}{p{0.23\linewidth}p{0.23\linewidth}p{0.22\linewidth}X}
    \toprule
    Evidence & Held fixed & Changed & Interpretation \\
    \midrule
    Pre/post grounding replay (136) & saved question, plan, KG, executor, oracle & deterministic grounding transformation & narrow runtime intervention; chiefly the additive-sum guard \\
    Checkpoint ladder (260) & questions, scorer, runtime budget, base within family & corpus, initial adapter, and objective across stages & sequential comparison, not matched ablation \\
    Hybrid/fully local runtime & questions and downstream runtime & Step-1 model and historical deployment configuration & system-configuration comparison; manifest limits causal attribution \\
    PACS A/B & program, oracle, and planner & surface rendering & paraphrase robustness \\
    WTQ base/v3/A/C & test set and official scorer & source and amount of supervision & supervision ladder, not single-component ablation \\
    \bottomrule
  \end{tabularx}
\end{table}

The existing artifacts therefore support two restrained conclusions.  First, deterministic
grounding can repair a specific class of saved sum plans at runtime.  Second, full SFT improves the
paired end-to-end checkpoint on both base models.  They do not identify the fraction of the SFT
gain caused by grounding-derived rows.  A causal claim would require corpora matched for unique
question, family, depth, output type, target length, and optimiser tokens while adding only ordinary
repair, grounding-derived repair, or grounding-derived preferences.  Other missing component
tests include guided versus unconstrained decoding, online grounding on/off, release checks and
feedback repair on/off, and deterministic fast path versus forced Step 2.  We list these as needed
experiments rather than retroactively labelling existing comparisons as ablations.

\section{Recommended Additional Evaluation}
\label{app:future-eval}

The matched grounding-data ablation should use at least three seeds and paired decoding on one
frozen test.  Before training, all four corpora should be matched by unique question, family,
program depth, output type, target length, and total optimiser tokens.  After training, each output
should be labelled for field, role, operator, return type, dependency, explicit-constraint recall,
hidden guard, validity, executability, answer correctness, and abstention.  Reporting a transition
matrix for each decision reveals whether an apparent accuracy gain came from the intended hard
signal or from generic JSON and domain-language adaptation.  Predicate-sensitive counterfactual
worlds should be part of the test so that dropping any evaluated constraint changes the denotation.
